% ============================================================
%  Inequações do 2º Grau — Apostila Premium | PratiqueJá
%  Engine: XeLaTeX
% ============================================================
\documentclass[12pt]{article}
\usepackage{../../pratiqueja}

\newcommand{\hlm}[1]{{\color{darkblue}#1}}

\setsubject{Inequações do 2.º Grau}

\begin{document}
\pagenumbering{arabic}
\setstretch{1.2}
\color{bodytext}

\heroheader{Inequações do 2.º Grau}%
           {Método gráfico, tabela de sinais e casos especiais}%
           {Matemática · Avançado}

\setlength{\columnsep}{18pt}
\setlength{\columnseprule}{0.5pt}
\def\columnseprulecolor{\color{exborder}}

\begin{multicols}{2}

\TW{Rev.}{badgepurple}{Revisão — Inequações do 1.º Grau}{Notação de intervalo e regra do negativo}

\regra{
Uma \hl{inequação} estabelece uma desigualdade com incógnita. Ao multiplicar ou dividir por um número \hl{negativo}, o sentido da desigualdade \textbf{inverte}.
}

\exemp{
$3x - 6 > 0 \;\Rightarrow\; x > 2 \;\Rightarrow\; \hlm{x \in (2, +\infty)}$

$-2x + 4 \leq 0 \;\Rightarrow\; x \geq 2 \;\Rightarrow\; \hlm{x \in [2, +\infty)}$
}

\nota{Notação: $(a,b)$ aberto, $[a,b]$ fechado, $(a,b]$ e $[a,b)$ semi-abertos.}

\TW{Graf}{darkblue}{Método Gráfico}{Usar a parábola para determinar o sinal}

\regra{
A inequação $ax^2 + bx + c \;\square\; 0$ pede os valores de $x$ onde a parábola $f(x) = ax^2 + bx + c$ está acima ($>0$), abaixo ($<0$) ou tocando ($=0$) o eixo $x$.

\begin{itemize}
  \item $a > 0$ (parábola $\cup$): $f(x) < 0$ \textbf{entre} as raízes; $f(x) > 0$ \textbf{fora}
  \item $a < 0$ (parábola $\cap$): $f(x) > 0$ \textbf{entre} as raízes; $f(x) < 0$ \textbf{fora}
\end{itemize}
}

\TW{1·2·3}{darkblue}{Método Passo a Passo}{Raízes → concavidade → solução}

\regra{
\textbf{1.} Passe tudo para um lado: $ax^2 + bx + c \;\square\; 0$

\textbf{2.} Calcule $\Delta$ e encontre as raízes $x_1 \leq x_2$ (se existirem)

\textbf{3.} Analise o sinal de $a$ para saber a concavidade

\textbf{4.} Monte a solução com base no sinal pedido
}

\exemp{
\textbf{Exemplo — $x^2 - 5x + 6 \leq 0$:}

$\Delta = 25 - 24 = 1 \;\Rightarrow\; x_1 = 2,\; x_2 = 3$

$a = 1 > 0$ → parábola $\cup$ → $f(x) \leq 0$ \emph{entre} as raízes

$\hlm{x \in [2,\, 3]}$

\textbf{Exemplo — $-x^2 + 4x - 3 > 0$:}

$\Delta = 16 - 12 = 4 \;\Rightarrow\; x_1 = 1,\; x_2 = 3$

$a = -1 < 0$ → parábola $\cap$ → $f(x) > 0$ \emph{entre} as raízes

$\hlm{x \in (1,\, 3)}$
}

\TW{Esp}{badgepurple}{Casos Especiais — $\Delta \leq 0$}{Solução é $\emptyset$ ou $\mathbb{R}$}

\regra{
Quando $\Delta < 0$, a parábola não cruza o eixo $x$. O sinal de $f(x)$ é constante (igual ao sinal de $a$) para todo $x$.

Quando $\Delta = 0$, existe apenas uma raiz (raiz dupla): $x_v$.
}

\exemp{
$x^2 + 1 > 0$: $\Delta < 0$; $a > 0$ → $f(x) > 0$ sempre

$\hlm{x \in \mathbb{R}}$

$x^2 + 1 < 0$: mesma função, sinal contrário ao pedido

$\hlm{S = \emptyset}$

$-x^2 + 4 \geq 0 \;\Leftrightarrow\; x^2 - 4 \leq 0$: $\Delta = 16$; $x_1 = -2$, $x_2 = 2$

$\hlm{x \in [-2,\, 2]}$

$(x-3)^2 \leq 0$: $\Delta = 0$, raiz dupla $x_v = 3$

$(x-3)^2 \geq 0$ sempre; $\leq 0$ só quando $(x-3)^2 = 0$

$\hlm{x = 3}$ (apenas o vértice)
}

\TW{( )( )}{badgepurple}{Inequação Fatorada}{Estudo de sinal por tabela}

\regra{
Quando $f(x)$ está na forma fatorada $a(x - x_1)(x - x_2)$, use uma \hl{tabela de sinais}: analise o sinal de cada fator separadamente e combine-os pelo produto.
}

\exemp{
\textbf{$(x-1)(x-4) \leq 0$ — zeros: $x = 1$ e $x = 4$:}

\vspace{4pt}
\noindent
\begin{minipage}[t]{\linewidth}
\footnotesize
\begin{tabular}{|c|c|c|c|c|c|}
\hline
$x$ & $({-}\infty,1)$ & $1$ & $(1,4)$ & $4$ & $(4,{+}\infty)$ \\
\hline
$(x-1)$ & $-$ & $0$ & $+$ & $+$ & $+$ \\
\hline
$(x-4)$ & $-$ & $-$ & $-$ & $0$ & $+$ \\
\hline
\textbf{Prod.} & $+$ & $0$ & \textcolor{green!50!black}{$-$} & $0$ & $+$ \\
\hline
\end{tabular}
\end{minipage}

\vspace{6pt}
$\hlm{\text{Solução }(\leq 0)\text{: } x \in [1,\, 4]}$
}

\end{multicols}

\end{document}
